\documentclass[11pt,a4paper]{article}
\usepackage[utf8]{inputenc}
\usepackage[T1]{fontenc}
\usepackage[english]{babel}
\usepackage{amsmath,amssymb,amsthm,mathtools}
\usepackage[margin=2.8cm]{geometry}
\usepackage{enumitem}
\usepackage{hyperref}
\hypersetup{colorlinks=true,linkcolor=black,citecolor=blue,urlcolor=blue}
\emergencystretch=2em
\theoremstyle{plain}
\newtheorem{theorem}{Theorem}[section]
\newtheorem{lemma}[theorem]{Lemma}
\newtheorem{proposition}[theorem]{Proposition}
\newtheorem{corollary}[theorem]{Corollary}
\theoremstyle{definition}
\newtheorem{definition}[theorem]{Definition}
\theoremstyle{remark}
\newtheorem{remark}[theorem]{Remark}
\newcommand{\RP}{\mathbb{RP}}
\newcommand{\PGL}{\mathrm{PGL}}
\newcommand{\rk}{\operatorname{rank}}
\newcommand{\tr}{\operatorname{tr}}
\title{Reflective projective billiards in higher dimensions:\\
monodromy, threshold phenomena, and explicit polytopes}
\author{Jorge Lucas Gonz\'alez\thanks{Email: jorge.lucas.glez@gmail.com}\\
\small Independent researcher}
\date{August 2026}
\begin{document}
\maketitle
\begin{abstract}
We develop a linear and algebro-geometric framework for reflective projective
billiards in projective dimension at least three.  Scalar products of harmonic
homologies lead to a parity obstruction, a sharp general lower bound for the
period, and a trichotomy between the rigid end, the threshold, and the stable
range.  We construct bounded rational polytopes with primitive open families of
periodic trajectories in dimensions three through six, analyze the invariant
quadrics and reflection groups of the examples, and reduce the threshold problem
in projective three-space to an explicit irreducibility question.  We also isolate
the determinantal equation behind the rank-three threshold argument and explain
exactly why it becomes vacuous in vector rank four.
\end{abstract}

\noindent\textbf{Keywords.} Projective billiards; harmonic homologies; reflection
groups; character varieties; periodic trajectories; convex polytopes.

\medskip
\noindent\textbf{2020 Mathematics Subject Classification.}
37C83, 51N15, 20F55, 14M12.

\section{Introduction}
Reflective projective billiards turn a dynamical closing condition into a
factorisation problem in \(\PGL_{n+1}\).  This paper studies the dimension-dependent
consequences of that observation.  The main themes are the number of equations in
the scalar-product locus, determinant parity, the rank of partial products, and
the geometry needed to convert a word of involutions into a bounded polytope with
a strict itinerary.

A projective billiard is \emph{$k$-reflective} when its billiard map has an open
set of $k$-periodic points; the first examples, of period three and of every even
period at least four, are due to Fierobe \cite{fie-ex}, and Problem~1 of
\cite{opb} asks which further periods occur.
Fierobe's higher-dimensional non-existence theorem for triangular families and
his subsequent analysis of caustics provide complementary background for the
higher-dimensional setting \cite{fie-tri,fie-caustics}.
The definitions and openness criterion are recalled below, and every construction
and auxiliary result used in the higher-dimensional arguments is proved here.
For related projective factorisation results, including products of three harmonic
homologies, see \cite{heidari-honari}.

The main contributions of this paper are: a dimension-dependent parity and rank
framework for scalar monodromy; explicit bounded rational polytopes with
primitive open families in dimensions three through six; exact analyses of their
invariant quadrics and reflection groups; and a reduction of the threshold
question in projective three-space to an explicit irreducibility problem.  The
irreducibility of $V_6$ and the resulting threshold theorem are kept explicitly
open where the proof does not close.
\section{Setting}\label{sec:setting}

A \emph{projective billiard} \cite{tab} is a bounded convex polytope
$\Omega\subset\mathbb{R}^n$ together with a \emph{transversal line field}
assigning to each point in the relative interior of a facet of $\partial\Omega$
a line $L(p)$ through $p$ not contained in the facet hyperplane; reflection is
defined by requiring the incoming ray, the outgoing ray, the facet tangent and
$L(p)$ to have cross ratio $-1$ in their impact pencil.

Throughout this paper we restrict to \emph{central facet fields}: on each
facet the transversal lines belong to the pencil through one fixed projective
centre outside the supporting hyperplane. No reduction from arbitrary
transversal fields to this class is asserted. Statements about the period
spectrum use the stated convention on the facets visited by the witness.

\begin{definition}[Fierobe \cite{fie-ex}]\label{def:krefl}
$\Omega$ is \emph{$k$-reflective} if there are a $k$-periodic orbit
$(p_1,\dots,p_k)$ and an \textbf{open} $U_1\times U_2\subset(\partial\Omega)^2$
containing $(p_1,p_2)$ every point of which begins a $k$-periodic orbit.
\end{definition}

The open set is the whole content of the definition, and it is where the difficulty
lives: exhibiting one periodic orbit is easy, exhibiting an open family of initial
impact pairs is not.

\subsection{Reflection as a harmonic homology}

For a covector $L$ and a vector $O$ with $LO\neq0$ put $H=I-2\,OL/(LO)$. Then
$H^2=I$, $HO=-O$, $Hx=x$ for $Lx=0$, and $\det H=-1$.

\begin{proposition}\label{prop:law}
Let $p$ lie on $L$, let $\tau\in\ker L\setminus\{0\}$ be a tangent direction, and
put $\nu=O-p$. The involution induced by $H$
on the pencil at $p$ sends $a\tau+b\nu$ to $a\tau-b\nu$; hence the cross ratio of
(incoming, outgoing; tangent, transversal) is $-1$. So reflection in the side with
transversal field the pencil through $O$ \emph{is} the harmonic homology with axis
the facet hyperplane and centre $O$.
\end{proposition}

\begin{proof}
$H$ fixes $p$ and is an involution, so induces an involution of the pencil at $p$
fixing the two distinct lines $\mathbb{R}\tau$ and $pO$; in the coordinate taking them to $0$
and $\infty$ it is $m\mapsto-m$. Parametrising by $(a:b)$ with direction
$a\tau+b\nu$, the tangent is $(1:0)$, the transversal $(0:1)$, and with
$[ij]=a_ib_j-a_jb_i$ the cross ratio of $(a:b),(a:-b),(1:0),(0:1)$ is
$[13][24]/([14][23])=-1$.
\end{proof}

Consequently, if $H_{j_k}\cdots H_{j_1}$ is scalar then \emph{every} line returns to
itself after the corresponding $k$ reflections. The converse is false: an isolated
trajectory may well close without the product being scalar, and it is precisely the
scalar case that produces an \emph{open set} of closed trajectories, which is what
Definition~\ref{def:krefl} demands. If the product is $cI$, determinants give
$c^{n+1}=(-1)^k$. For odd projective dimension $n$, the left side is positive,
so $k$ must be even. For even $n$, the real scalar is $c=(-1)^k$ and
determinants give no additional obstruction.
For the central fields considered here, an open family also forces scalar
monodromy. The initial lines of a regular family form an open subset of the
Grassmannian of projective lines. Their invariance under the fixed product
is an algebraic incidence condition; vanishing on an open subset forces
it on the entire Grassmannian. A linear map preserving every two-dimensional
vector subspace preserves every one-dimensional subspace, obtained as an
intersection of such subspaces in vector dimension at least three.
Applying this to basis vectors and their pairwise sums shows that the map
is scalar. This necessity concerns an open family; an isolated closed orbit
does not suffice.

\section{An openness criterion}

Everything below rests on this; it assumes no symmetry.

\begin{lemma}[Exit map]\label{lem:exit}
Let $\Omega$ be a convex polytope, $p$ in the relative interior of a facet, $w$ pointing into
$\Omega$. The ray $p+sw$, $s>0$, meets $\partial\Omega$ in exactly one point $E(p,w)$, and
$E$ is continuous wherever $E(p,w)$ lies in the relative interior of a facet.
\end{lemma}

\begin{proof}
$\Omega$ is convex and bounded, so $s^\ast=\sup\{s>0:p+sw\in\Omega\}$ is finite and
$E(p,w)=p+s^\ast w\in\partial\Omega$; uniqueness is convexity. If $E(p,w)$ is
in the relative interior of a facet $F$ then $w$ is not parallel to its
supporting hyperplane, so the intersection determinant is nonzero and Cramer's
rule gives $s^\ast$ near $(p,w)$ as a quotient of polynomials with nonvanishing
denominator. Since the relative interior of $F$ is open in $\partial\Omega$,
nearby data still exit through $F$.
\end{proof}

\begin{remark}
Requiring nonzero intersection determinants with every supporting hyperplane is unnecessarily restrictive: a strict trajectory may be parallel to facets it never reaches. Continuity only requires transversality at the actual exit facet, together with persistence of the strict inequalities for the other facets.
\end{remark}

\begin{theorem}[Openness criterion]\label{thm:open}
Let $\Omega$ be a convex polytope with facets $F_0,\dots,F_{k-1}$ and transversal
field on $F_i$ the pencil through $O_i$, with $L_iO_i\neq0$. Assume
\textup{(H1)} $H_{k-1}\cdots H_0=cI$, $c\neq0$; \textup{(H2)} there is an orbit
with itinerary $0,1,\dots,k-1$ whose impacts are strictly interior to their facets,
with strictly positive step parameters, each losing facet excluded by a strict
inequality and nonvanishing transversality at each impact; \textup{(H3)} the return
direction is $u_k=\kappa u_0$ with $\kappa>0$. Then $\Omega$ is $k$-reflective; if
the $k$ impacts lie on distinct facets the period is primitive.
\end{theorem}

\begin{proof}
Each step is a line--line intersection followed by a reflection, so the impact
parameters and directions are rational functions of the initial data, with
denominators nonvanishing at the witness by (H2); the vanishing locus being closed,
they are defined and continuous nearby. The finitely many strict inequalities of
(H2) persist, and by Lemma~\ref{lem:exit} each of the $k$ exit maps is continuous
with value in the relative interior of a facet, which is open in $\partial\Omega$; composing,
the itinerary is unchanged for nearby data. By (H1) the product is the identity of
$\PGL_{n+1}$, so \emph{every} line returns to itself: $\ell_k=\ell_0$ as an unoriented
line, and since $\ell_0$ is not the supporting hyperplane of $F_0$ the intersection
$\ell_0\cap F_0$ is a single point, whence $p_k=p_0$. For the orientation, (H1)
only gives $u_k=\kappa u_0$ with $\kappa$ real nonzero --- neither $\pm1$ nor a
priori locally constant --- but $\kappa=(u_k\cdot u_0)/(u_0\cdot u_0)$ is
continuous and nonvanishing and equals a positive number at the witness by (H3), so
$\kappa>0$ on a smaller neighbourhood. Finally, a proper period $d\mid k$, $0<d<k$,
would force $p_d=p_0$; but $p_d$ lies on facet $F_d$ and $p_0$ on $F_0$, which are
distinct facets, and the impacts are interior to them, so $p_d\neq p_0$.
\end{proof}

Hypotheses (H2)--(H3) are finitely many strict inequalities and (H1) one identity,
so they can be certified exactly for explicit data. We use this repeatedly
below.
\section{Higher dimensions}\label{sec:higher}

Nothing in the definitions or openness argument used $n=2$: a harmonic
homology of $\mathbb{RP}^n$ is $H=I-2OL/(LO)$ with $L$ a hyperplane covector, the
reflection law of Proposition~\ref{prop:law} is the same harmonic condition in the
pencil of lines at the impact point inside the $2$-plane spanned by the ray and the
transversal, and the openness criterion of Theorem~\ref{thm:open} goes through
verbatim for a convex polytope of $\mathbb{R}^n$ with impacts interior to facets:
its proof uses only convexity, Cramer's rule and finitely many strict
inequalities. What changes are the numbers, and they change in a way that explains
$k=4,5,6$ as one law.

\subsection{The four counts, and a trichotomy}

Write $N=\dim\PGL_{n+1}=(n+1)^2-1$. A harmonic homology of $\mathbb{RP}^n$ is a pair
$(L,O)$ with $LO\neq0$ modulo two scalings, so $\dim\mathrm{Inv}=2(n+1)-2=2n$.
At a point where the corresponding Jacobian ranks are maximal and the indicated
stabiliser is trivial, a direct dimension count gives:
\[
  \dim V_k=2nk-N,\qquad
  \dim(\text{moduli})=2nk-2N=2n(k-n-2),\qquad
  \dim S_k=nk+n .
\]
For the expected local value of $\dim S_k$: the space of quadrics of $\mathbb{RP}^n$ has dimension
$\binom{n+2}2-1$, each mirror is one $n_i\in\mathbb{RP}^n$, and the closure
condition is at most $\dim\mathrm{O}(Q)=n(n+1)/2$ conditions, so
$\dim S_k=\bigl[\binom{n+2}2-1\bigr]+nk-\frac{n(n+1)}2=nk+n$.

\begin{proposition}[Trichotomy of the expected counts]\label{prop:trichotomy}
Under the maximal-rank and stabiliser hypotheses just stated, the expected moduli
dimension is $0$ if and only if $k=n+2$, and the expected dimensions of $S_k$ and
$V_k$ agree if and only if $k=n+3$; for $k\ge n+4$, the expected dimension of
$S_k$ is smaller. These are the rigid, threshold, and dimension-deficient regimes,
respectively.
\end{proposition}

\begin{proof}
$2n(k-n-2)=0\iff k=n+2$. And $2nk-N=nk+n\iff nk=N+n=n^2+3n\iff k=n+3$; the
inequality for larger $k$ follows since the left side grows twice as fast.
\end{proof}

\subsection{Where the counts are valid, and by how much they fail}

The counts above are stated at points where the relevant Jacobian has maximal rank.
That hypothesis is not automatic, and it fails at three of the examples of these
notes. The following two statements say exactly when, and by exactly how much.

Write $\Phi$ for the map sending a $k$-tuple of harmonic homologies to its product in
$\PGL_m$, so that $V_k=\Phi^{-1}(\mathrm{id})$, and let $A_j=H_{k-1}\cdots H_{j+1}$
and $H'_j=A_jH_jA_j^{-1}$ be the \emph{unfolded} mirrors, with normals $a_j=A_jn_j$.

\begin{proposition}[The cokernel]\label{prop:cokernel}
Under the trace pairing,
\[
  \operatorname{coker}d\Phi\;\cong\;\{Y\in\mathfrak{sl}_m:\ [Y,H'_j]=0\ \text{for all }j\},
\]
the traceless part of the commutant of the group generated by the unfolded mirrors.
If moreover the mirrors are orthogonal for a positive definite form $Q$
and the $a_j$ span $\mathbb{R}^m$, then in coordinates with $Q=I$ this is
$\{S\ \text{symmetric traceless}:\ Sa_j\parallel a_j\ \forall j\}$, the space of
quadratic forms the configuration preserves beyond the given one.
\end{proposition}

\begin{proof}
Differentiating $H^2=I$ gives $\delta H\,H+H\,\delta H=0$, so $X=\delta H\,H$
satisfies $HXH=-X$; the space $W_H=\{X:HXH=-X\}$ is the $-1$ eigenspace of the
involution $\mathrm{Ad}_H$, of dimension $2(m-1)$, which is the number of parameters
of a harmonic homology, and $X\mapsto XH$ is bijective, so the tangent space maps
onto $W_H$. Varying only $H_j$,
$\delta P\cdot P^{-1}=A_j(\delta H_j H_j)A_j^{-1}=W_{H'_j}$, whence
$\operatorname{im}d\Phi=\sum_jW_{H'_j}$. Now $\mathrm{Ad}_{H'}$ is self-adjoint for
the trace form, so $W_{H'}^{\perp}$ is its $+1$ eigenspace, the centraliser of $H'$;
and the dual of $\mathfrak{pgl}_m$ under that pairing is $\mathfrak{sl}_m$. For the
last sentence, $[Y,I-2aa^{\top}/|a|^2]=0$ forces $Ya=\lambda a$ and
$Y^{\top}a=\lambda a$; splitting $Y$ into symmetric and antisymmetric parts and
subtracting gives $Ka=0$ for the antisymmetric part, hence $K=0$ when the $a_j$ span.
\end{proof}

\begin{proposition}[The corank is the number of blocks, minus one]\label{prop:corank}
Suppose the mirrors are orthogonal for a positive definite form, their
normals span $\mathbb R^m$, and their representation decomposes into $s$
pairwise non-isomorphic irreducible summands. Then $d\Phi$ has corank $s-1$,
and the tangent space to the scalar-product equations has dimension
\[
  \dim\ker d\Phi=2kn-\bigl((n+1)^2-1\bigr)+(s-1).
\]
For $s=1$ the map is a submersion and this is also the local dimension of
$V_k$. For $s>1$ it is only a tangent-space calculation; higher-order
equations may decrease the local dimension.
\end{proposition}

\begin{proof}
The unfolded mirrors generate the same group as the original ones, since
$H'_{k-1}=H_{k-1}$ and each $H_j$ is recovered inductively from the $H'$. By
Proposition~\ref{prop:cokernel} the cokernel is the traceless part of the commutant.
By the positive definite and spanning hypotheses in
Proposition~\ref{prop:cokernel}, every commuting endomorphism is symmetric.
The eigenspaces of a symmetric commuting endomorphism are invariant, so on
each irreducible block it is scalar. There are no off-diagonal intertwiners
between non-isomorphic blocks. The commutant therefore has dimension $s$,
and its traceless part has dimension $s-1$. Rank--nullity gives the formula;
the submersion theorem applies when $s=1$.
\end{proof}

\begin{proposition}[A quadratic obstruction at the rank-five example]
\label{prop:quadratic-corank}
At the configuration of Theorem~\ref{thm:rp4b5}, the tangent space to the
scalar-product equations has dimension $33$, whereas the complex local
dimension is $32$. The real local dimension is at most $32$.
\end{proposition}

\begin{proof}
Use independent affine coordinates for the seven poles $p_i$ and axes $l_i$,
fixing the first nonzero coordinate of each at its value in the displayed
root $n_i$. There are $7\cdot8=56$ coordinates. For
$H_i=I-2p_il_i/(l_ip_i)$ and $P=H_6\cdots H_0$, exact differentiation at
$p_i=n_i$, $l_i=n_i^\top$ gives $P=-I$ and $\operatorname{rank}dP=23$.
The matrix $Y=\operatorname{diag}(1,1,1,-4,1)$ commutes with every mirror,
and $\operatorname{tr}(Y\,dP)=0$ for all $56$ coordinate directions.

Here is an explicit tangent direction, with indices starting at zero; all
unlisted coordinate derivatives are zero:
\[
\begin{array}{c|rrrrrrr}
\text{coordinate}&p_{0,3}&l_{0,3}&l_{1,3}&p_{2,1}&p_{2,4}&l_{2,1}&l_{2,2}\\
\hline
\text{derivative}&2&2&2&-1&1&-1&1
\end{array}
\]
Along the straight coordinate path $tv$, rational multiplication gives
\[
 P'(0)=0,\qquad
 \operatorname{tr}(Y(P(t)+I))
 =\frac{2t^2(18t^2-5)}{(t^2+1)(2t^2+1)}=-10t^2+O(t^4).
\]
The supplement's verifier \texttt{check\_corank\_obstruction.py}
reproduces these identities and the rank in exact arithmetic.

Eliminate $23$ independent equations by the analytic implicit function
theorem. Their zero set is smooth of dimension $33$ and has tangent space
$\ker dP$. A curve in that set with initial derivative $v$ may change the
second derivative of $P$ by an element of $\operatorname{im}dP$ only.
Pairing with $Y$ annihilates that change. Thus the remaining scalar equation
has a nonzero quadratic term, and the local zero set has dimension at most
$32$. Finally $P$ takes values in the smooth determinant-$-1$ hypersurface
of dimension $24$, so its fibre is locally defined by $24$ equations in
$56$ variables. The complex Krull bound is at least $32$ on every component.
The complex local dimension is consequently exactly $32$.
\end{proof}

\begin{remark}[Which of our examples are which]\label{rem:whichblocks}
Of the five higher-dimensional configurations certified in this paper, two are
irreducible and three are not; the split is exactly the one predicted by
Theorem~\ref{thm:bmstops}: the
reducible ones are precisely the words of $B_m$ with isolated coordinates.
\[
\begin{array}{lccc}
\text{configuration} & s & \text{corank} & \text{counts valid?}\\\hline
\mathbb{R}^3,\ k=6\text{ (Thm.~\ref{thm:rp3example})} & 1 & 0 & \text{yes}\\
\mathbb{R}^3,\ k=6\text{ (Thm.~\ref{thm:rp3d4})} & 1 & 0 & \text{yes}\\
\mathbb{R}^4,\ k=7\text{ (Thm.~\ref{thm:rp4b5})} & 2 & 1 & \text{no}\\
\mathbb{R}^5,\ k=8\text{ (Thm.~\ref{thm:rp5b6})} & 3 & 2 & \text{no}\\
\mathbb{R}^6,\ k=9\text{ (Thm.~\ref{thm:rp6b7})} & 4 & 3 & \text{no}
\end{array}
\]
The block structure is the same in the last three: one $D_4$ block of rank four, plus
one line for each isolated coordinate. The three lines carry pairwise different
characters of the group --- on each, one generator acts by $-1$ and the others
trivially --- so they are non-isomorphic and
Proposition~\ref{prop:corank} applies. Checked over $\mathbb{Q}$ by
\path{verificadores/bloques.py} and \path{verificadores/tolerancias.py}.

None of this disturbs the examples themselves: they are certified by the openness
criterion, which uses no dimension count. What it disturbs is the right to apply the
counts, and Proposition~\ref{prop:trichotomy} with them, at those points.
\end{remark}

\subsection{Parity: the obstruction reverses}

When $n$ is even, $-I$ is a scalar in odd vector rank, so odd $k$ carries no
determinant obstruction. In odd projective dimensions the parity restriction is
different, as the theorem shows.

\begin{theorem}[Parity]\label{thm:parity}
A harmonic homology of $\mathbb{RP}^n$ has a canonical lift
$H=I-2OL/(LO)$, which fixes its axis pointwise and satisfies $H^2=I$ and
$\det H=-1$. When $n+1$ is even, this determinant is independent of the
choice between the two involutive lifts $H$ and $-H$. If the canonical lifts satisfy
$H_{k-1}\cdots H_0=cI$ then
\[
  c^{\,n+1}=(-1)^k .
\]
Hence for $n$ even there is no parity constraint on $k$ --- which is why $\mathbb{RP}^2$ admits
odd $k$ --- while for \textbf{$n$ odd, $k$ must be even}. In particular in
$\mathbb{RP}^3$ there is no $k$-reflective projective billiard with $k$ odd; the
rigid case $k=n+2=5$ is \emph{empty}, not merely degenerate.
\end{theorem}

\begin{proof}
Take determinants. For $n+1$ even, $\det(-H)=(-1)^{n+1}\det H=\det H$, so the sign
of the lift does not matter. If $n+1$ is even then $c^{n+1}>0$ for real $c\ne0$,
forcing $(-1)^k>0$.
\end{proof}

\begin{remark}[The quiver picture, and where the parity actually sits]
\label{rem:parityquiver}
It is worth setting Theorem~\ref{thm:parity} inside the Deligne--Simpson
machinery, because the same set-up is used in the next paragraph. Rescale all the
factors by one $\mu$ so that the product is the identity of
$\mathrm{GL}_{n+1}$: the local classes become $\mu$ with multiplicity $n$ and
$-\mu$ with multiplicity $1$, the associated star-shaped quiver of
\cite{cb-icm} has central dimension $n+1$ and $k$ legs of dimension
$\operatorname{rank}(A_i-\xi_{i1})=1$, so $\alpha=(n+1;1,\dots,1)$, and
Crawley-Boevey's necessary condition reads
\[
  \xi[\alpha]=\prod_{i=1}^k\prod_{j}\xi_{ij}^{\,m_{i,j-1}-m_{ij}}
  =\prod_{i=1}^k\mu^{\,n}\,(-\mu)
  =(-1)^k\mu^{\,k(n+1)}=1 .
\]
This is precisely $\prod_i\det A_i=1$; over $\mathbb C$ it can always be
arranged by the choice of $\mu$, so it is not an obstruction there. The parity
is a statement about the \emph{real} problem: $\mu$ must be real, and when $n+1$
is even $\mu^{k(n+1)}>0$, forcing $(-1)^k=1$. So the quiver normalisation
recovers Theorem~\ref{thm:parity}, but by the same determinant identity rather
than by an independent route.

The same computation verifies the other two conditions of \cite[Thm.~10]{cb-icm} for
$n=2$, $k=5$: $\alpha=(3;1^5)$ is a positive imaginary root, and
$p(\alpha)=2$ exceeds the largest value $1$ attained by
$p(\beta)+p(\gamma)+\cdots$ over the fifteen nontrivial decompositions of $\alpha$
into positive roots with $\xi=1$. So there exists a solution that is
\emph{irreducible as a representation}: five harmonic homologies with scalar product
whose mirrors have no common invariant subspace, for every $k$ satisfying the parity
condition, with no construction needed. That is not the same as irreducibility of the
\emph{variety} $V_k$, which is what Theorem~\ref{thm:reduction} needs, and the
criterion says nothing about it.
\end{remark}

\subsection{The rigid end, in every dimension}

The low-dimensional case needed in the argument is the following elementary fact.

\begin{lemma}[Four reflections in dimension three]\label{lem:fourref}
Let $R_0,R_1,R_2,R_3$ be orthogonal reflections of a positive definite
three-dimensional space and suppose $R_0R_1R_2R_3$ is scalar. Then their normals
span a subspace of dimension at most two.
\end{lemma}

\begin{proof}
The product has determinant $+1$, so the scalar is $I$. Thus
$R_0R_1=R_3R_2$. If this product is the identity, each pair consists of the same
reflection and the four normals span at most a two-dimensional space. Otherwise
the common product is a nonidentity rotation with a unique fixed axis. The axis
of a product of two reflections is perpendicular to both normals, so all four
normals lie in the perpendicular plane.
\end{proof}

\begin{theorem}[$k=n+2$ is always degenerate]\label{thm:rigidend}
Let $Q$ be positive definite on $\mathbb{R}^{m}$, $m=n+1\ge3$, let
$R_0,\dots,R_{k-1}$ be $Q$-reflections with $k=n+2=m+1$, and suppose
$R_0R_1\cdots R_{k-1}=cI$. Then $c=1$, $m$ is odd, and the normals
$n_0,\dots,n_{k-1}$ do \textbf{not} span $\mathbb{R}^m$. Consequently all $k$
mirrors contain a common projective point, and the region they bound is unbounded:
there is no such billiard.
\end{theorem}

\begin{proof}
The product lies in $\mathrm{O}(Q)$, and $cI\in\mathrm{O}(Q)$ forces $c=\pm1$.
Taking determinants, $c^m=(-1)^k=(-1)^{m+1}$; with $c=-1$ this reads
$(-1)^m=(-1)^{m+1}$, impossible, so $c=1$ and $m$ is odd. (For $m$ even the
statement is vacuous by Theorem~\ref{thm:parity}.)

Suppose the normals span $\mathbb{R}^m$. Put
\[
  V=\langle n_0,n_1,n_2\rangle,\quad W=\langle n_3,\dots,n_{k-1}\rangle,\quad
  X=R_0R_1R_2,\quad Y=R_3\cdots R_{k-1},
\]
so $XY=I$ and $X=Y^{-1}$; note $k-3=m-2\ge1$, so $W$ is defined. Each $R_i$ with
$i\le2$ fixes $n_i^\perp$ pointwise, so $X$ fixes $V^\perp$ pointwise; likewise
$Y^{-1}$ fixes $W^\perp$ pointwise. Hence $X$ fixes
$V^\perp+W^\perp=(V\cap W)^\perp$ pointwise. Since $V+W=\mathbb{R}^m$,
\[
  \dim(V\cap W)=\dim V+\dim W-m\le 3+(m-2)-m=1 ,
\]
so $\operatorname{codim}\mathrm{Fix}(X)\le1$. As $\det X=(-1)^3=-1$ we have
$X\neq I$, so $\mathrm{Fix}(X)$ is exactly a hyperplane; an element of
$\mathrm{O}(Q)$ fixing a hyperplane pointwise and not the identity acts by $-1$ on
the orthogonal line, so $X=R'$ is a \emph{reflection}, with normal $n'$ spanning
$\mathrm{Fix}(X)^\perp$. From $\mathrm{Fix}(X)\supseteq(V\cap W)^\perp$ we get
$n'\in V\cap W$, so $\dim(V\cap W)=1$ and the inequality above is an equality:
$\dim V=3$.

Now $R_0R_1R_2R'=I$ with $n_0,n_1,n_2,n'$ all in $V$, which is $3$-dimensional and
 on which $Q$ is definite; each of the four maps preserves $V$ and restricts to a
 reflection of it. By Lemma~\ref{lem:fourref} the four normals are coplanar, so
$\dim V\le2$ --- contradicting $\dim V=3$.

Finally, if the normals span a proper subspace $U$ then every mirror contains
$\mathbb{P}(U^\perp)\neq\emptyset$. Pick a point $p$ in it: if $p$ is affine the
region is a cone with apex $p$, and if $p$ is at infinity with direction $d$ every
mirror is invariant under translation by $d$ and so is the region. Either way it is
unbounded.
\end{proof}

\begin{remark}[Arbitrary signature]\label{rem:rigidsig}
Definiteness of $Q$ was used in exactly two places, and both can be replaced by an
explicit hypothesis, so the theorem holds over any nondegenerate $Q$ --- in
particular over $\mathbb{C}$, where all nondegenerate forms are equivalent:
\begin{enumerate}[label=\textup{(\alph*)},nosep]
  \item $\mathbb{R}^m=V\oplus V^\perp$ with each $R_i$ preserving both, which needs
        $V=\langle n_0,n_1,n_2\rangle$ \textbf{nondegenerate};
  \item ``an element of $\mathrm{O}(Q)$ fixing a hyperplane pointwise and not the
        identity acts by $-1$ on the orthogonal line'', which needs
        $\mathrm{Fix}(X)$ \textbf{nondegenerate}, i.e.\ its orthogonal line
        non-isotropic.
\end{enumerate}
Everything else is signature-free: the identity
$(U_1\cap U_2)^\perp=U_1^\perp+U_2^\perp$ holds for any nondegenerate $Q$, and
$\mathrm{Fix}(R_1R_0)=n_0^\perp\cap n_1^\perp$ exactly whenever $n_0,n_1$ are
independent, since $R_1R_0x=x$ forces
$\tfrac{\langle x,n_0\rangle}{\langle n_0,n_0\rangle}n_0
=\tfrac{\langle x,n_1\rangle}{\langle n_1,n_1\rangle}n_1$ and hence both
coefficients to vanish --- no appeal to signature anywhere. Moreover the relation is
invariant under cyclic rotation of the mirrors, since rotating conjugates the
product, so the split may be taken at any of the $k$ positions and it suffices that
\emph{some} three cyclically consecutive normals span a nondegenerate $3$-space.
Whether all $k$ rotations can fail simultaneously we do not determine. The billiard
application below uses a positive definite $Q$, so these signature issues do not
arise.
\end{remark}

\begin{remark}
Only the case $j=3$ of the splitting is used, and that is what reduces the argument
to Lemma~\ref{lem:fourref}: three reflections have determinant $-1$, which
is what rules out $X=I$ and forces $\mathrm{Fix}(X)$ to be a genuine hyperplane. A
split into two even halves would allow $X=I$ and prove nothing.
\end{remark}

\subsection{The threshold: what is proved, and what is not}

At $k=n+3$ the expected dimensions coincide. The following explicit example
illustrates the threshold case in dimension three.

\begin{proposition}[$\mathbb{RP}^3$, $k=6$]\label{prop:rp3}
There is an explicit configuration of six $Q$-reflections of $\mathbb{R}^4$ with
integer normals,
\[
\begin{gathered}
  n_0=(1,0,0,-3),\quad n_1=(2,1,-1,1),\quad n_2=(1,3,2,0),\\
  n_3=(0,-1,4,0),\quad n_4=(0,-55,50,17),\quad n_5=(-171,215,-40,-82),
\end{gathered}
\]
with $Q=x_0^2+x_1^2+x_2^2+x_3^2$ and $R_5R_4R_3R_2R_1R_0=-I$ exactly over
$\mathbb{Q}$, whose normals have rank $4$, at which
\[
  \dim V_6=6k-15=21,\qquad \dim S_6=3k+3=21,\qquad \dim(\text{moduli})=6k-30=6 ,
\]
these being the values on the component through the point. Consequently $S_6$
contains the irreducible component of $V_6$ through that point: \textbf{every
configuration in that component preserves a quadric}. The invariant form is
nondegenerate on a nonempty open subset of this component, since it is positive
definite at the displayed point.
\end{proposition}

\begin{proof}
The three Jacobian ranks and the rank of the linear system for the invariant quadric
are computed over $\mathbb{Q}$ by Gaussian elimination at the displayed point:
Jacobian of the $15$ equations has rank $15$, the infinitesimal stabiliser is
spanned by $I$, the Jacobian in the $n_i$ directions has rank
$6=\dim\mathrm{O}(Q)$, and the $24\times10$ system for $Q$ has rank $9$, so the
invariant quadric is unique up to scale. Semicontinuity propagates each value to the
component. Since $\dim S_6=\dim V_6$ and $S_6\subseteq V_6$ is closed, $S_6$ is a
union of components of $V_6$, and it contains the given point.
\end{proof}

\subsection{A $6$-reflective projective billiard in $\mathbb{R}^3$}

The counts and Proposition~\ref{prop:rp3} say where to look. Here is an actual
example. Its priority relative to the literature is not asserted here.

\begin{theorem}\label{thm:rp3example}
Let $Q=x_0^2+x_1^2+x_2^2+x_3^2$, let $n_i=(a_i,d_i)$ be the six integer vectors
\[
\begin{gathered}
 n_0=(-1,3,-4,-1),\quad n_1=(-1,0,4,1),\quad n_2=(3,4,-1,-2),\\
 n_3=(0,-8,91,0),\quad n_4=(0,2198,2119,8345),\quad n_5=(-1577,276,1117,-146),
\end{gathered}
\]
let $\varepsilon=(+,+,+,+,-,+)$, and let
\[
  \Omega=\{x\in\mathbb{R}^3:\ \varepsilon_i(a_i\cdot x+d_i)<0,\ i=0,\dots,5\},
\]
with the transversal field on facet $i$ the pencil of lines through the pole
$O_i=n_i$ of that facet. Then $\Omega$ is a \textbf{bounded} convex polytope and
the projective billiard on it is \textbf{$6$-reflective}, with primitive period
$6$.
\end{theorem}

\begin{proof}
The hypotheses of Theorem~\ref{thm:open}, whose proof applies verbatim to a convex
polytope of $\mathbb{R}^n$ with impacts interior to facets, are verified in exact
rational arithmetic.

\emph{\textup{(H1)}.} Each $R_i=I-2n_in_i^\top/(n_i\!\cdot\!n_i)$ satisfies
$R_i^2=I$, $\tr R_i=2$ and $R_i^\top R_i=I$, and
$R_0R_1R_2R_3R_4R_5=-I$ over $\mathbb{Q}$, computed in both orders. The six normals
have rank $4$ and the six facets are pairwise distinct.

\emph{Boundedness.} The recession cone is $\{y:\varepsilon_ia_i\cdot y\le0\ \forall
i\}$, and it is $\{0\}$ exactly when the inward normals positively span
$\mathbb{R}^3$, i.e.\ when they have rank $3$ and some $\mu\in\mathbb{Q}^6_{>0}$
satisfies $\sum_i\mu_i\varepsilon_ia_i=0$. Both hold, with
$\mu=(\tfrac{8056}{65},\tfrac{5581}5,\tfrac{61038}{65},1,2,1)$ and the identity
checked exactly. Nonemptiness: $x_0=(-\tfrac{49}{100},-\tfrac{153}{100},
-\tfrac{121}{100})$ satisfies all six inequalities strictly.

\emph{\textup{(H2)}.} The witness orbit starts at
$p_0=(-\tfrac12,-\tfrac32,-\tfrac54)$, which lies on facet $0$ and strictly inside
the other five, with direction $u_0=(-\tfrac49,-\tfrac7{10},\tfrac59)$ pointing
inwards. Its itinerary is $5,4,3,2,1,0$; at each of the six steps the first hit
(least step parameter) is the predicted facet, with $s>0$ strictly, the impact
strictly interior to that facet, every losing facet excluded by a strict
inequality, no non-winning facet parallel to the ray, and nonvanishing
transversality.

\emph{\textup{(H3)}.} After six steps $p_6=p_0$ and $u_6=u_0$ exactly, so
$\kappa=1>0$. The six impacts lie on six distinct facets, so the period is
primitive.

Reflection is implemented projectively: at an impact $q$ on facet $i$ the outgoing
direction is $w-w_4q$ where $(w,w_4)=R_i\,(u,0)^\top$, which is the harmonic
condition of Proposition~\ref{prop:law} in the $2$-plane spanned by the ray and the
transversal.
\end{proof}

\begin{proposition}[The explicit $\mathbb{R}^3$ billiard is not isolated]
\label{prop:rp3localfamily}
In the space $V_6$ of six harmonic homologies with scalar product, the configuration
of Theorem~\ref{thm:rp3example} is a smooth real point of dimension $21$. Moreover,
there is a relative open neighbourhood of it in $V_6$ all of whose configurations
define bounded convex $6$-reflective projective billiards with the same primitive
itinerary. Its local moduli space modulo $\PGL_4$ has dimension $6$.
\end{proposition}

\begin{proof}
Use independent affine coordinates for the axes and centres of the six homologies.
The exact Jacobian of the $15$ traceless scalar-product equations at the displayed
rational configuration has rank $15$; hence the implicit-function theorem gives a
smooth local $36-15=21$ dimensional locus. The infinitesimal stabiliser has
dimension $1$ in $\mathfrak{gl}_4$, consisting only of scalar matrices, so the
projective orbit has dimension $15$ and the local quotient has dimension $6$.

All the remaining conditions in the proof of Theorem~\ref{thm:rp3example} are
strict: a positive dependence among the inward normals, a strictly interior point,
strict first-hit inequalities, interior impacts, nonvanishing transversality, and
positive return multiplier. They persist under sufficiently small perturbations.
The product relation supplies the closing line, and the same continuity argument as
in Theorem~\ref{thm:open} supplies the oriented return. Thus the openness criterion
applies throughout a smaller neighbourhood. The rank computation uses the same
exact rational Jacobian calculation as the verifier for Proposition~\ref{prop:rp3},
now evaluated at the billiard configuration itself.
\end{proof}

\begin{remark}
This lands exactly at the threshold: $k=6=n+3$ for $n=3$. And $k$ is even, as
Theorem~\ref{thm:parity} requires for odd $n$ --- there is no $5$-reflective
example in $\mathbb{R}^3$ to be found, which is why the rigid case is empty rather
than merely degenerate. Two things are \emph{not} claimed. This configuration and
the one of Proposition~\ref{prop:rp3} are different, and we do not claim they lie on
the same component of $V_6$; and irreducibility of $V_6$ in dimension~$3$ is open,
as is the threshold theorem itself (Corollary~\ref{cor:rp3reduction}). What is proved
here is existence.

Finally, this example is not a kaleidoscope in disguise: its normals are rational,
so each $\operatorname{tr}(R_aR_b)=2+2\cos\phi_{ab}$ is rational, and by Niven's theorem all
fifteen values of $\cos\phi_{ab}$ fall outside $\{0,\pm\frac12,\pm1\}$. So every
one of the fifteen pairs generates an infinite dihedral group and the mirror group
is infinite.
\end{remark}

The configuration is produced by a factorisation argument with the splitting
$6=3+3$ and an elementary property of $\mathrm{O}(4)$:

\begin{lemma}\label{lem:o4}
Every element of $\mathrm{O}(4)$ of determinant $-1$ has eigenvalue $1$, hence is a
product of at most three reflections.
\end{lemma}

\begin{proof}
In the block normal form the number of $(-1)$-blocks is odd, so $1$ or $3$; in both
cases $4=2a+\#\{\pm1\text{ blocks}\}$ forces at least one $(+1)$-block.
\end{proof}

Choose rational normals $n_0,n_1,n_2$ and set
$X=-(R_2R_1R_0)^{-1}\in\mathrm O(4)$, so $\det X=-1$. By Lemma~\ref{lem:o4},
$X$ fixes a nonzero rational vector. Choose a rational $v$ with $Xv\ne v$ and put
$u=Xv-v$. The reflection $R_u$ interchanges $v$ and $Xv$, so
$Y=R_uX\in\mathrm{SO}(4)$ fixes $v$ and restricts to an element of
$\mathrm{SO}(v^\perp)\cong\mathrm{SO}(3)$. If this restriction is not the
identity, let $e$ span its fixed axis and choose a rational
$a\in v^\perp\cap e^\perp$. The reflection $R_a$ fixes $e$ and reverses the
orientation of $e^\perp$, so $R_aYR_a=Y^{-1}$. Put $R_b:=YR_a$. Then $R_b^2=I$,
it fixes $e$, and its determinant on $v^\perp$ is $-1$, so its other eigenvalues
are $+1$ and $-1$; hence $R_b$ is a reflection and $Y=R_bR_a$. If the restriction is
the identity, choose any rational reflection $R_a$ of $v^\perp$ and put
$R_b=R_a$. Thus $Y=R_bR_a$, and
\[
              R_uR_bR_aR_2R_1R_0=-I.
\]
All choices and operations are rational, giving the claimed construction.

\subsection{The threshold reduces to irreducibility}

The following dimension argument isolates the additional hypothesis needed to turn
the threshold count into a theorem: irreducibility of the scalar-product locus.

\begin{theorem}[The threshold reduces to irreducibility]\label{thm:reduction}
Fix $n\ge2$ and $k=n+3$, and suppose there is a point of $V_k\subset\mathbb{RP}^n$
at which
\begin{enumerate}[label=\textup{(\alph*)},nosep]
  \item the Jacobian of the $N=(n+1)^2-1$ scalar-product equations has rank $N$, so
        that the component through it has dimension $2nk-N$; and
  \item the component of the locus $S_k$ of configurations preserving a nonzero
        quadratic form through this point has dimension $nk+n$.
\end{enumerate}
If $V_k$ is irreducible, then
\[
  V_{n+3}=S_{n+3} :
\]
\textbf{every} configuration of $n+3$ harmonic homologies of $\PGL_{n+1}$ with
scalar product preserves a nonzero quadric.
\end{theorem}

\begin{proof}
By (a) and irreducibility, $\dim V_k=2nk-N$. By Proposition~\ref{prop:trichotomy},
$k=n+3$ gives $nk+n=2nk-N$, so by (b) $\dim S_k\ge\dim V_k$. But $S_k$ is a closed
subvariety of the irreducible $V_k$, and a closed subvariety of an irreducible
variety of the same dimension is the whole variety.
\end{proof}

For $n=3$ the hypotheses (a) and (b) are the content of
Proposition~\ref{prop:rp3}, certified in exact arithmetic at the displayed
rational point. So:

\begin{corollary}\label{cor:rp3reduction}
If $V_6\subset\mathbb{RP}^3$ is irreducible, then every configuration of six
harmonic homologies of $\PGL_4$ with scalar product preserves a quadric --- the
threshold theorem in dimension three.
\end{corollary}

This is worth stating because it moves the problem to where there are tools. By the
dictionary of Section~\ref{sec:higher}, irreducibility of $V_k$ is irreducibility of
a relative character variety of the $k$-punctured sphere with involution monodromies,
which is the multiplicative Deligne--Simpson problem for the star-shaped quiver with
central dimension $n+1$ and $k$ legs of dimension $1$.

\subsection{Why the grouping argument could not have generalised}

The pair grouping gives a useful dimension count, but it cannot by itself settle the
threshold in higher dimensions. The pair lemma is as follows.

\begin{lemma}[Pair lemma in dimension $m$]\label{lem:pairgen}
Let $H_0,H_1$ be harmonic homologies of $\mathbb{RP}^{m-1}$ with axes $L_0\neq L_1$
and let $A=H_1H_0$. Then $A$ fixes $L_0\cap L_1$ pointwise and, if $A$ is regular
in the sense that its eigenvalues are $\lambda,\lambda^{-1},1,\dots,1$ with
$\lambda\neq\pm1$ and the $1$-eigenspace is exactly $L_0\cap L_1$, the quadrics
preserved by $A$ are exactly those preserved by \emph{both} $H_0$ and $H_1$, and
they form a linear space of dimension
\[
  1+\binom{m-1}2 .
\]
A single harmonic homology preserves a linear space of quadrics of dimension
$1+\binom m2$.
\end{lemma}

\begin{proof}
Both $H_i$ fix $L_0\cap L_1$ pointwise, hence so does $A$; and
$H_0AH_0=A^{-1}$, so $H_0$ interchanges the $\lambda$- and
$\lambda^{-1}$-eigenlines and preserves the $1$-eigenspace. In a basis
$(e_\lambda,e_{\lambda^{-1}},f_3,\dots,f_m)$ adapted to this, $\tau_A(Q)$ multiplies
$Q_{ij}$ by $\lambda_i\lambda_j$, so $Q$ is $A$-invariant exactly when $Q_{ij}=0$
unless $\lambda_i\lambda_j=1$: this leaves $Q_{12}$ and the whole block on the
$1$-eigenspace, of dimension $1+\binom{m-1}2$. In the same basis
$H_0=\left(\begin{smallmatrix}0&q&\\r&0&\\&&I\end{smallmatrix}\right)$ with $qr=1$,
since $H_0^2=I$ and $H_0$ is the identity on $L_0\supseteq L_0\cap L_1$. For
$Q=2\beta x_1x_2+Q'$ with $Q'$ on the $1$-eigenspace, $Q(H_0x)=2\beta qr\,x_1x_2+Q'=Q$.
So every $A$-invariant quadric is $H_0$-invariant, hence also $H_1=AH_0$-invariant;
the converse is immediate. For a single homology, in a basis where
$H=\operatorname{diag}(-1,1,\dots,1)$ the invariant quadrics are those with
$Q_{1j}=0$ for $j\ge2$: dimension $1+\binom m2$.
\end{proof}

\begin{proposition}[The deficit]\label{prop:deficit}
Inside the projective space of quadrics of $\mathbb{RP}^{m-1}$, of dimension
$\binom{m+1}2-1$, the locus $\mathcal Q_A$ of Lemma~\ref{lem:pairgen} has
codimension $2m-2$. Hence $p$ pairs of mirrors can be expected to share a quadric
only when
\[
  p\,(2m-2)\le\binom{m+1}2-1,\qquad\text{i.e.}\qquad p\le\frac{m+2}4 .
\]
At the threshold $k=n+3=m+2$ the grouping uses $p=\lfloor(m+2)/2\rfloor$ pairs, and
the shortfall $p(2m-2)-\bigl(\binom{m+1}2-1\bigr)$ is
\[
\begin{array}{c|cccccc}
 m       & 3 & 4 & 5 & 6 & 7 & 8\\\hline
 \text{deficit} & 3 & 9 & 10 & 20 & 21 & 35
\end{array}
\]
\end{proposition}

\begin{proof}
$\binom{m+1}2-1-\binom{m-1}2=2m-2$, and the rest is arithmetic.
\end{proof}

\begin{remark}
At $m=3$ the tabulated deficit is $3$, but the grouping there is $2+2+1$, not
$2+2$. The fifth homology is $-N^{-1}$, so the real requirement is that one pencil
$\mathcal P_A$ meet the polarity condition for $N$, which is $m-1=2$ conditions on a
$1$-parameter pencil: a shortfall of exactly $\mathbf1$. And the relation pays
exactly one condition, namely $\det(B-I)=0$, which holds automatically because
$B=H_3H_2$ is a product of two involutions. That coincidence --- one condition
needed, one condition supplied --- is a special feature of $m=3$. For $m\ge4$ the shortfall is $9$ or more while the
relation still supplies one condition. So the grouping argument does not merely
resist generalisation: it cannot generalise. Theorem~\ref{thm:reduction} is the way
through, and it needs irreducibility instead.
\end{remark}

\begin{remark}[What irreducibility of $V_6$ would take]\label{rem:v6irr}
One can try a fibration over the first three factors. Split $6=3+3$: the relation reads
$H_5H_4H_3=c(H_2H_1H_0)^{-1}=:X$, and $H_5X=H_4H_3$ must be a product of \emph{two}
harmonic homologies, which by Lemma~\ref{lem:pairgen} means
$\operatorname{rank}(H_5X-I)\le2$.

Now comes the point. Writing $H_5=I-2w\tilde w/s$,
\[
  H_5X-I=(X-I)-\tfrac2s\,w\,(\tilde wX),
\]
a \textbf{rank-one perturbation} of $X-I$. A rank-one perturbation lowers rank by at
most one, so $\operatorname{rank}(H_5X-I)\le2$ \emph{forces}
$\operatorname{rank}(X-I)\le3$, i.e.\ $\det(X-I)=0$: $X$ must have eigenvalue $1$.
For general harmonic homologies that is \textbf{not} automatic --- a product of three
of them has determinant $-1$ but need not fix a point --- and it is one condition.
Given it, $\operatorname{rank}(X-I)=3$ and the three further conditions for a
rank-one perturbation of a rank-$3$ matrix to have rank $2$ are that $w$ lie in
$\operatorname{im}(X-I)$, that $(\tilde wX)^\top$ lie in
$\operatorname{im}\bigl((X-I)^\top\bigr)$, and one scalar normalisation. So the count
is
\[
\begin{aligned}
  &\underbrace{18-1}_{H_0,H_1,H_2\ \text{with}\ \det(X-I)=0}
  +\underbrace{6-3}_{H_5\ \text{in the determinantal locus}}
  +\underbrace{1}_{\text{factorising}\ H_4H_3}\\
  &\hspace{4cm}=17+3+1=21=\dim V_6 .
\end{aligned}
\]

And the asymmetry that appears here is exactly the reason $S_6$ can be a
\emph{component} rather than a thin subvariety. On $S_6$ the mirrors are
$Q$-reflections, so $X\in\mathrm{O}(Q)$ with $\det X=-1$, and then
Lemma~\ref{lem:o4} makes $\det(X-I)=0$ \textbf{automatic}: the $3+3$ split is free
there and costs nothing, while off $S_6$ it costs a condition. That is where the
excess dimension comes from.

What remains, and it is now a question about two explicit varieties rather than
about billiards: is the hypersurface $\{\det(X-I)=0\}\subset\mathrm{Inv}^3$
irreducible, and is the determinantal locus of admissible $H_5$ irreducible of
dimension $3$ over it? With both, $V_6$ is irreducible and
Corollary~\ref{cor:rp3reduction} gives the threshold theorem in dimension three.

Numerical evidence points that way: Newton's method from sixteen random starts
converged to sixteen points of $V_6$, and \emph{every one} of them preserved a
quadric, with the smallest singular value of the linear system at the level of
machine epsilon while the next one was of order $10^{-1}$. That is evidence, not
proof, and it is recorded as such.
\end{remark}

\subsection{A second example in $\mathbb{R}^3$, and the dichotomy repeats}

The example of Theorem~\ref{thm:rp3example} is not a kaleidoscope. There is also
one with data in $\{0,\pm1\}$, whose mirror group is finite.

\begin{theorem}\label{thm:rp3d4}
Let $Q=x_0^2+x_1^2+x_2^2+x_3^2$ and take the six mirrors with normals
\[
\begin{gathered}
 n_0=(1,0,1,0),\quad n_1=(1,0,-1,0),\quad n_2=(1,0,0,-1),\\
 n_3=(0,1,1,0),\quad n_4=(0,1,-1,0),\quad n_5=(0,1,0,-1),
\end{gathered}
\]
which are six of the twelve mirrors of the reflection group $D_4$, with
$\varepsilon=(-,-,+,-,-,+)$ and $\Omega=\{x:\varepsilon_i(a_i\cdot x+d_i)<0\}$. Then
$\Omega$ is a bounded convex polytope and the projective billiard on it is
$6$-reflective with primitive period $6$; the witness orbit starts at
$p_0=(\tfrac47,\tfrac45,-\tfrac47)$ with direction $u_0=(15,5,3)$, has itinerary
$2,1,3,5,4,0$, and closes with $p_6=p_0$ and $\kappa=1$ exactly over $\mathbb{Q}$.
Its mirror group is \textbf{finite}, of order $96$ in $\PGL_4$ and $192$ in
$\mathrm{O}(4)$: it is a spherical kaleidoscope of $\mathbb{RP}^3$.
\end{theorem}

\begin{proof}
As in Theorem~\ref{thm:rp3example}, all hypotheses of Theorem~\ref{thm:open} are
certified in exact rational arithmetic: the six maps are $Q$-reflections whose
product in the order of the itinerary is $\pm I$; the recession cone is $\{0\}$,
by rank $3$ of the inward normals together with an explicit
$\mu\in\mathbb{Q}^6_{>0}$ with $\sum\mu_i\varepsilon_ia_i=0$; the point
$(\tfrac12,\tfrac12,-\tfrac14)$ is interior; and at each of the six steps the first
hit is the predicted facet, strictly interior to it, with every losing facet excluded
strictly. Finiteness of the group is a breadth-first closure in $\PGL_4$, exact.
\end{proof}

\begin{remark}
In dimension three these examples exhibit both behaviours: Theorem~\ref{thm:rp3d4}
has a finite reflection group, while Theorem~\ref{thm:rp3example} has an infinite
mirror group. The latter shows that the moduli space is not a list of symmetries.

A remark on how it was found, because it is a methodological point. A first sweep over
the mirrors of $D_4$, $B_4$ and $F_4$ found nothing, because it fixed
$\varepsilon_0=+1$ to halve the work --- and $\varepsilon\mapsto-\varepsilon$ is not a
symmetry of the problem, it exchanges a polytope for a different region. With all
$2^k$ sign patterns, and comparing the itineraries the dynamics actually produces
against \emph{all} orderings of the chosen mirrors whose product is scalar rather
than only the cyclic ones, the example appears at once.
\end{remark}

\subsection{A $7$-reflective projective billiard in $\mathbb{R}^4$}

The threshold $k=n+3$ is also attained one dimension higher.  The following
example is deliberately given with rational data, including a rational witness.

\begin{theorem}\label{thm:rp4b5}
Let $Q_0=X_0^2+\cdots+X_4^2$ and let
\[
\begin{gathered}
 n_0=(0,1,0,0,-1),\ n_1=(0,0,1,0,-1),\ n_2=(0,0,0,1,0),\\
 n_3=(1,-1,0,0,0),\ n_4=(0,1,1,0,0),\ n_5=(1,0,0,0,1),\\
 n_6=(1,0,-1,0,0).
\end{gathered}
\]
These are roots of $B_5$.  Put $R_i=I-2n_i n_i^T/(n_i^Tn_i)$, let
\[
 T=\begin{pmatrix}
 -1/3&-1/3&1/3&1/3&1/3\\
 1&0&0&0&0\\0&1&0&0&0\\0&0&1&0&0\\0&0&0&1&0
 \end{pmatrix},\qquad Q=T^T T,
\]
and use the affine chart whose last homogeneous coordinate is $1$.  The seven
walls, oriented inward, are
\[
\begin{array}{rclcrcl}
 -x+w&<&0,&&-y+w&<&0,\\
 z&<&0,&&4x+y-z-w-1&<&0,\\
 -x-y&<&0,&&x+y-z-4w-1&<&0,\\
 &&& &x+4y-z-w-1&<&0.
\end{array}
\]
Then their projective billiard, with the polar fields for $Q$, is a bounded
convex $7$-reflective projective billiard of primitive period $7$.  A rational
witness starts on the fifth wall at
\[
 p_0=(0,0,-4/7,-2/21),\qquad
 u_0=(1/5,1,-1/4,-1/5),
\]
and has successive first impacts on the walls
\[
 5,6,0,1,2,3,4.
\]
It closes with $p_7=p_0$ and $u_7=u_0$.
\end{theorem}

\begin{proof}
The explicit $B_5$ calculation gives
\[
 R_6R_5R_4R_3R_2R_1R_0=-I_5.
\]
Thus $H_i=T^{-1}R_iT$ are $Q$-reflections and their product is projectively
scalar.  The displayed affine inequalities are precisely the oriented images of
the $n_i^\perp$, cleared of denominators so as to be legible.  The positive
dependence is in the \emph{unscaled} normalisation: taking the covectors
$n_i^\top T$ with the sign that makes the interior negative, their seven spatial
parts sum to zero, so $\mu=(1,\dots,1)$ is a strictly positive dependence and
the recession cone is trivial; those seven vectors have rank four.  Clearing
denominators rescales three of the seven and destroys the all-ones dependence,
so the certificate must be read before that step.  Finally
$(1/10,1/10,-1/10,0)$ is strictly interior.  So the domain is nonempty, bounded
and convex.

For the stated $(p_0,u_0)$, exact rational ray tracing gives the seven positive
flight times
\[
 \frac4{189},\ \frac{74}{867},\ \frac{1050}{10693},\ \frac{105}{703},\
 \frac{1575}{49799},\ \frac{17340}{112703},\ \frac{20}{129}.
\]
At each step the indicated wall is the unique first wall and the impact is
strictly interior relative to it.  Applying $H_i$ to the homogeneous tangent
vector and returning to the affine chart gives $p_7=p_0$ and $u_7=u_0$ exactly.
The strict inequalities persist under perturbation of the initial data, so
Theorem~\ref{thm:open} gives an open set of $7$-periodic orbits.  Since the
seven walls are distinct, the exhibited period is primitive.
\end{proof}

\begin{remark}[Reproducibility; calculation added by Codex, 8 August 2026]
Every equality and strict inequality in Theorem~\ref{thm:rp4b5} is checked over
$\mathbb{Q}$ by \path{verificadores/reflectivo_rp4_b5.py}.  In particular,
the theorem does not rely on the floating-point search that located the word.
\end{remark}

\subsection{An $8$-reflective projective billiard in $\mathbb{R}^5$}

The preceding word extends once more, and in this case the extension is also
realised dynamically.

\begin{theorem}\label{thm:rp5b6}
There is a bounded convex $8$-reflective projective billiard in $\mathbb{R}^5$
with primitive period $8$, preserving a positive definite quadric.  It may be
chosen with rational data.  More precisely, take in $\mathbb{R}^6$ the roots
\[
\begin{gathered}
(0,1,0,0,-1,0),(0,0,1,0,-1,0),(0,0,0,1,0,0),(1,-1,0,0,0,0),\\
(0,1,1,0,0,0),(1,0,0,0,1,0),(1,0,-1,0,0,0),(0,0,0,0,0,1),
\end{gathered}
\]
with orientations $(+,+,-,+,-,-,+,-)$ and the affine chart defined by
$\ell=(-1,1,1,1,3,1)$.  The associated $Q$-reflections have product $-I_6$.
The point and direction
\[
p_0=(-2/33,0,0,4/11,4/11),\qquad
u_0=(-11/10,1,-2/5,7/10,1/5)
\]
start on wall $4$ and have first-impact itinerary
\[
5,6,7,0,1,2,3,4,
\]
closing exactly with $p_8=p_0$ and $u_8=u_0$.
\end{theorem}

\begin{proof}
The first seven roots are those of Theorem~\ref{thm:rp4b5}, extended by zero,
and their product is $\operatorname{diag}(-I_5,1)$.  The final coordinate
reflection makes the product $-I_6$.  In the stated chart the eight inward
spatial normals have positive dependence with all coefficients $1$, and the
point $(-1/10,1/10,1/10,1/10,1/10)$ is strictly interior, proving boundedness
and nonemptiness.  Exact rational tracing from the displayed data has the stated
itinerary, with every competing impact strictly later and every impact interior
relative to its wall; it returns exactly to $(p_0,u_0)$.  Theorem~\ref{thm:open}
then supplies an open set of $8$-periodic trajectories.
\end{proof}

\begin{remark}[Reproducibility; calculation added by Codex, 8 August 2026]
The full list of affine inequalities, flight times and strict comparisons for
Theorem~\ref{thm:rp5b6} is checked over $\mathbb{Q}$ by
\path{verificadores/reflectivo_rp5_b6.py}.
\end{remark}

\subsection{A $9$-reflective projective billiard in $\mathbb{R}^6$}

The threshold is attained once more, two dimensions higher, and with $k$ odd.

\begin{theorem}\label{thm:rp6b7}
Let $Q=\sum_{j=1}^{7}x_j^2$ and let the nine mirrors of $\mathbb{RP}^6$ have normals
\[
\begin{aligned}
  n_0&=(0,0,0,0,0,1,-1), & n_1&=(0,1,0,1,0,0,0), & n_2&=(0,0,1,0,0,0,0),\\
  n_3&=(0,0,0,1,0,1,0),  & n_4&=(1,0,0,0,0,0,0), & n_5&=(0,1,0,0,0,0,-1),\\
  n_6&=(0,0,0,0,1,0,0),  & n_7&=(0,0,0,1,0,0,-1),& n_8&=(0,1,0,0,0,1,0),
\end{aligned}
\]
with transversal field on face $i$ the pencil through the pole $O_i=n_i$ of $Q$, and
take the affine chart $\ell\cdot X=1$ with $\ell=(-1,1,-1,-3,1,1,-1)$, solving for the
fourth homogeneous coordinate, and the sign vector
$\varepsilon=(-,+,+,+,+,-,-,+,-)$. The resulting bounded convex polytope of
$\mathbb{R}^6$ is a $9$-reflective projective billiard.
\end{theorem}

\begin{proof}
The nine reflections are involutions of $\mathrm{O}(Q)$ with $R_8\cdots R_0=-I_7$,
which is (H1). The polytope is bounded because the nine outward normals in the chart
\emph{sum to zero} --- so $\mu=(1,\dots,1)$ is a strictly positive certificate --- and
have rank six. The point
\[
p_0=(-\tfrac6{19},\tfrac7{228},-\tfrac6{19},0,
\tfrac7{228},-\tfrac7{228})
\]
lies on face $6$ and strictly inside the other eight. Exact rational tracing from
$p_0$ in the direction
\[
u_0=10^{-6}(-1003425,-270995,31421,1824019,-819948,-545098)
\]
gives the itinerary
$7,8,0,1,2,3,4,5,6$, with every impact strictly interior to its face, every competing
face excluded by a strict inequality, all flight parameters positive, and a return to
$p_0$ with $\kappa=1$. Theorem~\ref{thm:open} applies.
\end{proof}

\begin{remark}
Three points. First, $k=9$ is odd and $n=6$ is even, which is exactly what
Theorem~\ref{thm:parity} permits; in odd ambient dimension $k$ would have been forced
even. Second, $9=n+3$, so this is again the threshold, as in
Theorems~\ref{thm:rp3example}, \ref{thm:rp3d4}, \ref{thm:rp3d4b}, \ref{thm:rp4b5} and
\ref{thm:rp5b6}. Third, and this is the informative part, the word obeys
Theorem~\ref{thm:bmstops} on the nose: six of its roots are of type $e_i\pm e_j$ and
live on the coordinates $\{1,3,5,6\}$, forming a $D_4$ block of rank four, and the
other three are the sign changes $e_0,e_2,e_4$. The block has the same multiset of
$\cos^2$ of pairwise angles as the cores of Theorems~\ref{thm:rp4b5} and
\ref{thm:rp5b6} and as the configuration of Theorem~\ref{thm:rp3d4b}: it is the same
rank-four object again, and the three isolated coordinates make the configuration
reducible, of corank three. Checked over $\mathbb{Q}$ by
\path{verificadores/reflectivo_rp6_b7.py}.
\end{remark}

\subsection{Why the $B_m$ route stops at rank four}

The words used in Theorems~\ref{thm:rp4b5} and \ref{thm:rp5b6} live in $B_5$ and
$B_6$, and it is natural to ask how far that supply goes. Not far, and the reason is
a counting argument inside $S_m$.

\begin{theorem}[The $B_m$ threshold words are reducible for $m\ge5$]\label{thm:bmstops}
Let $m\ge5$ and let $r_1\cdots r_{m+2}=-I_m$ be a word of $m+2$ reflections of
$B_m$. Then the configuration is reducible: some proper nonzero subspace of
$\mathbb{R}^m$ is invariant under all the $r_i$. The same holds in $D_m$.
\end{theorem}

\begin{proof}
Write $B_m=(\mathbb{Z}/2)^m\rtimes S_m$. A reflection in $e_i$ has trivial
permutation part; a reflection in $e_i\pm e_j$ has permutation part the transposition
$(ij)$. Since $-I_m$ has trivial permutation part, the $t$ transpositions occurring
must multiply to the identity of $S_m$.

If some coordinate $c$ occurs in no root of type $e_i\pm e_j$ then $\langle
e_c\rangle$ is invariant under every $r_i$; and if the graph whose edges are those
transpositions is disconnected, the span of the coordinates of one component is
invariant. So irreducibility forces that graph to be connected and to span all $m$
vertices, that is, the transpositions act transitively.

A transitive factorisation of the identity of $S_m$ into $t$ transpositions satisfies
$t=2(m-1)+2g$ with $g\ge0$, hence $t\ge2(m-1)$; and $t\le m+2$. Therefore
$2(m-1)\le m+2$, that is $m\le4$. In $D_m$ every root is of type $e_i\pm e_j$, so
$t=m+2$ and the same inequality applies.
\end{proof}

\begin{corollary}\label{cor:onlyE7E8}
Let $m\ge5$ and let a word of $m+2$ reflections with product $-I_m$ generate a
finite reflection group acting irreducibly on $\mathbb{R}^m$. Then that group is $E_7$
(with $m=7$) or $E_8$ (with $m=8$).
\end{corollary}

\begin{proof}
The group has rank $m$, acts irreducibly and contains $-I_m$, which is the product.
Among the irreducible finite reflection groups containing $-I$, those of rank at
least five are $B_m$, $D_m$ with $m$ even, $E_7$ and $E_8$;
Theorem~\ref{thm:bmstops} removes the first two.
\end{proof}

The inequality is sharp and says more than the statement. At $m=4$ it forces $t=6$
and no sign changes at all, so every irreducible threshold word of rank four consists
of six roots of type $e_i\pm e_j$: a $D_4$ word. Deleting the isolated coordinates
from the words of Theorems~\ref{thm:rp4b5} and \ref{thm:rp5b6} leaves exactly such a
core, and it is the same core in both. The uniform supply of $m+2$ reflections of
$B_m$ with product $-I_m$, for every $m\ge5$, is therefore that one rank-four
configuration padded with coordinate reflections, and its algebraic content does not
grow with $m$. That the core is more than algebra is the next example.

\begin{theorem}\label{thm:rp3d4b}
The bounded convex polytope of $\mathbb{R}^3$ cut out by the six faces with normals
\[
  n_0=(0,1,0,-1),\quad n_1=(0,0,1,-1),\quad n_2=(1,-1,0,0),
\]
\[
  n_3=(0,1,1,0),\quad n_4=(1,0,0,1),\quad n_5=(1,0,-1,0)
\]
in the affine chart $\ell=(-2,-2,-2,-1)$, with transversal field on face $i$ the
pencil through the pole $O_i=n_i$ of $Q=\sum_j x_j^2$, is a $6$-reflective projective
billiard. It is not projectively equivalent to the billiard of
Theorem~\ref{thm:rp3d4}.
\end{theorem}

\begin{proof}
The six reflections satisfy $R_5\cdots R_0=-I_4$, which is (H1). The polytope is
bounded, by an explicit positive rational combination of the outward normals summing
to zero together with those normals having rank three, and nonempty, the point
$(-\tfrac19,-\tfrac19,\tfrac19)$ being strictly interior to all six half-spaces. The
ray from that point in direction $(\tfrac54,-\tfrac14,2)$ meets face $3$ first;
starting there, exact rational tracing gives the itinerary $4,5,0,1,2,3$ with every
impact strictly interior to its face, every competing face excluded by a strict
inequality, all flight parameters positive, and a return to the starting point with
$\kappa=1$. Theorem~\ref{thm:open} applies. For the last claim, both configurations
are irreducible, so each preserves a unique quadric up to scale; a projective
equivalence between them would carry $Q$ to a multiple of $Q$ and hence be orthogonal
up to scale, but the multiset of $\cos^2$ of the angles between normals has three
zero entries here and six for Theorem~\ref{thm:rp3d4}.
\end{proof}

\begin{remark}
Its mirror group is all of $W(D_4)$, of order $192$, exactly as for
Theorem~\ref{thm:rp3d4}. So this is a second kaleidoscope in $\mathbb{R}^3$ rather
than a third kind of example, and like the symmetric pentagon of the plane it unfolds
to a finite group and has \emph{every} orbit periodic. The example with infinite
mirror group in $\mathbb{R}^3$ remains the one of Theorem~\ref{thm:rp3example}.
Checked over $\mathbb{Q}$ by \path{verificadores/reflectivo_rp3_d4b.py}.
\end{remark}

\subsection{A rank bound, and a lower bound on \texorpdfstring{$k$}{k}}

Before returning to irreducibility, one lemma that costs a line and constrains
everything.

\begin{lemma}[Rank bound]\label{lem:rankbound}
For the canonical lifts $H_i=I-2O_iL_i/(L_iO_i)$ of harmonic homologies of $\PGL_m$,
\[
  \rk\bigl(H_{j-1}\cdots H_0-I\bigr)\le j .
\]
\end{lemma}

\begin{proof}
The product telescopes:
$H_{j-1}\cdots H_0-I=\sum_{i=0}^{j-1}H_{j-1}\cdots H_{i+1}(H_i-I)$, a sum of $j$
matrices each of rank $1$, since $H_i-I=-2O_iL_i/(L_iO_i)$.
\end{proof}

\begin{corollary}\label{cor:rankbound}
\begin{enumerate}[label=\textup{(\alph*)},nosep]
  \item A product of $j<m$ harmonic homologies has eigenvalue $1$ with multiplicity
        at least $m-j$. In particular a product of three of them in $\PGL_4$ always
        has eigenvalue $1$.
  \item If $H_{k-1}\cdots H_0=cI$ with $c\neq1$, then $\rk(cI-I)=m\le k$, i.e.
        \[
          \boxed{\;k\ge n+1\;}
        \]
        There are no $k$-reflective projective billiards of $\mathbb{RP}^n$ with
        $k<n+1$ and nontrivial scalar. For $n=2$ this is $k\ge3$, which is sharp:
        Fierobe's triangle.
  \item Consequently, in the splitting $6=3+3$ of Remark~\ref{rem:v6irr} the
        condition $\det(X-I)=0$ is one genuine condition when $c=-1$ --- eigenvalue
        $1$ comes free by \textup{(a)}, and eigenvalue $-1$ does not --- and is
        automatic when $c=+1$.
\end{enumerate}
\end{corollary}

\begin{proof}
(a) is Lemma~\ref{lem:rankbound} read as $\dim\ker(M-I)\ge m-j$. (b) is the same
lemma applied to the product, since $\rk(cI-I)=m$ for $c\neq1$. (c) follows from (a)
with $j=3$, $m=4$, plus the computation
$\det(X-I)=\pm\det(M-cI)$ for $X=cM^{-1}$.
\end{proof}

\subsection{A conditional reduction for the regular part of $V_6$}

Use determinant-one matrix lifts and work over $\mathbb C$. Put
\[
 \mathcal P=\{M\in\mathrm{SL}_4:\operatorname{rank}(M-I)\le2\}.
\]
Every product of two normalized harmonic homologies belongs to $\mathcal P$.
The distinction between the full determinantal locus and the image of the
factorisation map matters at exceptional matrices.

\begin{lemma}[The pair locus and its regular fibres]\label{lem:Pirr}
$\mathcal P$ is irreducible of dimension $11$. On its nonempty open subset
$\mathcal P^{\mathrm{rs}}$ of matrices with eigenvalues
$1,1,\lambda,\lambda^{-1}$, $\lambda\ne0,\pm1$, and two-dimensional
$1$-eigenspace, the factorisation map from pairs of harmonic homologies
is smooth with geometrically irreducible one-dimensional fibres.
The fibre over $I$ instead has dimension $6$.
\end{lemma}

\begin{proof}
Consider the incidence space of pairs $(F,M)$ with
$F\in\mathrm{Gr}(2,4)$ and $M|_F=I$. In local coordinates for the Grassmannian,
such a matrix has blocks $\left(\begin{smallmatrix}P&0\\R&I\end{smallmatrix}\right)$,
with $P\in\mathrm{SL}_2$ and arbitrary $R$. This is a locally trivial
irreducible family of dimension $4+3+4=11$, surjecting onto $\mathcal P$.
Over the rank-two open subset $F=\ker(M-I)$ is unique, proving the dimension.

For $M\in\mathcal P^{\mathrm{rs}}$, write $F=\ker(M-I)$ and let $E$ be
the sum of the $\lambda$ and $\lambda^{-1}$ eigenlines. A factor $H_0$ obeys
$H_0MH_0=M^{-1}$, so preserves $F,E$ and interchanges those eigenlines.
Its restriction to $E$ already has a one-dimensional negative eigenspace;
since a harmonic homology has exactly one negative eigenline, it is identity
on $F$. In an eigenbasis of $E$ it is
$\left(\begin{smallmatrix}0&q\\q^{-1}&0\end{smallmatrix}\right)$, $q\ne0$.
Conversely each such choice gives the pair $H_1=MH_0,H_0$.
The fibre is $\mathbb G_m$. Ordering the simple eigenvalues and choosing
eigenvectors etale locally proves smoothness. For $M=I$ one simply has
$H_1=H_0$, and all homologies are allowed, a six-dimensional family.
\end{proof}

Group the six mirrors into $A=H_1H_0$, $B=H_3H_2$, $C=H_5H_4$.
Scalar monodromy forces $C=c(BA)^{-1}$ and
$\operatorname{rank}(BA-cI)\le2$. Define
\[
 \mathcal B_c=\{(A,B)\in\mathcal P^2:
                   \operatorname{rank}(BA-cI)\le2\},
\]
and let $\mathcal B_c^{\mathrm{rs}}$ impose in addition
$A,B,c(BA)^{-1}\in\mathcal P^{\mathrm{rs}}$.
Above this open subset the six-mirror locus is a smooth family of
geometrically irreducible fibres of dimension $3$.

\begin{proposition}[Conditional global reduction]\label{prop:v6reduce}
Suppose $\mathcal B_c^{\mathrm{rs}}$ is nonempty and irreducible and every
component of $V_6$ meets its inverse image. Then $V_6$ is irreducible;
under the hypotheses of Corollary~\ref{cor:rp3reduction} this implies its
threshold conclusion. Neither hypothesis on the global components is
established here.
\end{proposition}

\begin{proof}
The smooth factorisation map over the indicated irreducible base has
geometrically irreducible fibres. Smoothness makes it open, so all components
dominate the base, and irreducibility of the geometric generic fibre leaves
exactly one. By the assumed intersection with every component, this open
subset is dense in all of $V_6$; its closure is therefore all of $V_6$.
\end{proof}

The expression $11+11-4+3=21$ is an expected dimension count, not a proof
that the pulled-back rank condition has codimension four. The dimension
$21$ is known at the submersive examples by the independent differential
calculation. Exceptional pair fibres and possible additional components
must be controlled before promoting the regular reduction to a global one.

Finally, a local parameterisation of harmonic homologies will be useful.

\begin{lemma}[Rank-one chart for harmonic homologies]\label{lem:invchart}
For a real vector space $V$ of dimension $m$, every projective harmonic homology
has a lift
\[
             N=2w\tilde w/(\tilde w w)-I,
\]
where $w\in V$, $\tilde w\in V^*$ and $\tilde w w\ne0$. This identifies the
space of such homologies with the open subset $\tilde w w\ne0$ of
$\mathbb P(V)\times\mathbb P(V^*)$, which is irreducible of dimension
$2m-2$.
\end{lemma}

\begin{proof}
The displayed map fixes $w$ and acts by $-I$ on $\ker\tilde w$, so it represents
a homology; changing its sign gives the usual involutive lift that fixes its
axis. Conversely, the centre line and axis hyperplane of any harmonic homology
give $w$ and $\ker\tilde w$, and their non-incidence is exactly $\tilde w w\ne0$.
The parameter space is a nonempty open subset of a product of projective spaces.
\end{proof}

\begin{proposition}[The determinantal locus is a $(1,1)$-hypersurface]\label{prop:11hyp}
Let $Y\in\PGL_4$ with $\rk(Y-I)=3$, write $C=Y-I$, let $\kappa$ span $\ker C$ and
$\ell$ span $\ker C^\top$. In the parameterisation of
Lemma~\ref{lem:invchart}, the locus
$\mathcal D_Y=\{H\in\mathrm{Inv}:\rk(HY-I)\le2\}$ is cut out by
\[
  \textup{(i)}\ \ell w=0,\qquad
  \textup{(ii)}\ \tilde w\kappa=0,\qquad
  \textup{(iii)}\ 2\,b(\tilde w)^\top w=\tilde ww,
\]
where $b(\tilde w)$ solves $b^\top C=\tilde wY$ and depends linearly on $\tilde w$.
Conditions \textup{(i)} and \textup{(ii)} are linear, so $\mathcal D_Y$ is a
\textbf{bidegree $(1,1)$ hypersurface in $\RP^2\times\RP^2$}, irreducible of
dimension $3$ whenever its coefficient matrix has rank $\ge2$.
\end{proposition}

\begin{proof}
Writing $H=I-2w\tilde w/s$ with $s=\tilde ww$, we get
$HY-I=C-\tfrac2s\,w\,(\tilde wY)$, a rank-one perturbation of $C$. By Wedderburn's
rank-one reduction, if $u=Ca$ and $v^\top=b^\top C$ then $\rk(C-t\,uv^\top)=\rk C-1$
exactly when $t\,b^\top Ca=1$; and the perturbation lowers the rank at all only if
$u\in\operatorname{im}C$ and $v^\top\in\operatorname{im}C^\top$. The first is
$\ell w=0$. For the second, $\tilde wY$ must annihilate $\ker C$, and
$Y\kappa=\kappa$ gives $\tilde wY\kappa=\tilde w\kappa$, which is (ii). The scalar
condition is $\tfrac2s\,b^\top w=1$, i.e.\ (iii), and $b^\top w$ is well defined
because $b$ is determined modulo $\ell$ and $\ell w=0$. Bidegree $(1,1)$ forms are
reducible only if they factor as a form in $w$ times a form in $\tilde w$, i.e.\ only
in rank $1$.
\end{proof}

\begin{remark}
At the configuration of Theorem~\ref{thm:rp3example} this is checked in exact
arithmetic: $\rk(X-I)=3$, both kernels are lines, and the $3\times3$ coefficient
matrix of \textup{(iii)} is
\[
  \begin{pmatrix}0&-\tfrac{65}{24}&-\tfrac{11}{24}\\
  \tfrac{65}{24}&0&-\tfrac7{12}\\ \tfrac{11}{24}&\tfrac7{12}&0\end{pmatrix},
\]
antisymmetric of rank $2$: so $\mathcal D_X$ is irreducible there. What is left for
the threshold theorem in dimension three includes irreducibility of the
regular base and control of the exceptional components, as made explicit
in Proposition~\ref{prop:v6reduce}.
\end{remark}

\subsection{Which values of $k$ admit a metric example}

Existence and \emph{metric} existence are different questions, and it is the second
one that has structure. Call a $k$-reflective projective billiard \emph{metric} if it
preserves a definite conic; by Proposition~\ref{prop:law} that is the same as being
the metric billiard of the corresponding round projective metric; its
spherical lift is round by the positive definite Cayley--Klein construction.

Three facts, of which the first two are new here and the third is
Theorem~\ref{thm:rigidend} read correctly.

\begin{proposition}[$k=m$: metric, and unique]\label{prop:kmunique}
Let $k=m=n+1$, suppose the configuration is metric, and suppose its lifted
product is $-I_m$. Then the normals are
pairwise orthogonal, so the configuration is the orthant bounded by $m$ mutually
orthogonal mirrors, unique up to isometry. It is realised for every $n$.
\end{proposition}

\begin{proof}
Conjugate the definite form to the Euclidean inner product and argue by
induction on $m$, with $m=1$ immediate. Multiplication of the product relation
on the right by $H_0$ gives $H_{m-1}\cdots H_1=-H_0$. Telescoping shows
\[
 \operatorname{im}(-H_0-I)\subset
 \operatorname{span}(n_1,\ldots,n_{m-1}).
\]
For this inclusion, use
$A_r\cdots A_1-I=(A_r-I)A_{r-1}\cdots A_1+(A_{r-1}\cdots A_1-I)$,
so every summand has image in one of the displayed normal lines.
The left side equals $n_0^\perp$ and has dimension $m-1$, so the two spaces
are equal. Thus every other normal is orthogonal to $n_0$. Restriction to
$n_0^\perp$ gives a product of $m-1$ reflections equal to $-I$ there;
the induction hypothesis finishes the proof. For the
realisation take $H_i=I-2e_ie_i^{\ast}$ on the simplex $\{x_i>0\}$, with transversal
field on the facet $x_i=0$ the pencil through the opposite vertex $[e_i]$: the
homologies commute, their product is $-I$, and the line
$f(s,t)=(s-\lambda_0t,\dots,s-\lambda_nt)$ with $\lambda_0>\dots>\lambda_n>0$ crosses
the facets one at a time, in strict order, with all other coordinates on the correct
side of the orthant after folding.
\end{proof}

For $n=2$ this is the classification of $3$-reflective spherical billiards by three
orthogonal great circles \cite{bknz}; the proposition is its extension to every
dimension, and the point is not the construction but the uniqueness, which comes for
free from saturation of the rank bound.

\begin{proposition}[Common centre kills every nondegenerate conic]\label{prop:commoncentre}
A configuration of harmonic homologies with a common centre $O$ and at least two
distinct axes preserves no nondegenerate quadric.
\end{proposition}

\begin{proof}
If $S$ is invariant and nondegenerate, each $H_i$ is an involution preserving $S$, so
its $+1$ and $-1$ eigenspaces are $S$-orthogonal. The $-1$ eigenspace is
$\langle O\rangle$ and the $+1$ eigenspace is $\ker L_i$, so $SO\in\langle
L_i\rangle$. As $S$ is nondegenerate and $O\neq0$ we have $SO\neq0$, hence
$L_i\propto SO$ for every $i$ and all the axes coincide.
\end{proof}

\begin{corollary}[The rigid end is not empty; it is empty of metric billiards]
\label{cor:rigidnotempty}
Theorem~\ref{thm:rigidend} says there is no metric $k$-reflective billiard with
$k=m+1$. It does not say the value $k=m+1$ is unattainable: centrally projective
configurations, all of whose centres coincide, realise it, and by
Proposition~\ref{prop:commoncentre} they preserve no nondegenerate conic at all. So
at the rigid end the absence of an invariant conic is not a choice of construction:
it is forced.
\end{corollary}

\begin{remark}[The metric table]\label{rem:metrictable}
Collecting, with $m=n+1$:
\[
\begin{array}{lll}
 k=m   & \text{metric with product }-I\text{, and then unique} &
 \text{Proposition~\ref{prop:kmunique}}\\
 k=m+1 & \text{never metric} & \text{Theorem~\ref{thm:rigidend}, Corollary~\ref{cor:rigidnotempty}}\\
 k=m+2 & \text{metric examples exist} & \text{Section~\ref{sec:higher}}\\
 k\ge m+3 & \text{not classified here} & \text{}
\end{array}
\]
The determinant condition for a metric configuration with product $-I$ is
$k\equiv m\pmod 2$. Beyond the values proved in this paper, the existence of
metric examples is left open here. The displayed finite computations are checked
over $\mathbb{Q}$ by
\path{verificadores/simplex_ortogonal.py} and
\path{verificadores/extremo_rigido_no_metrico.py}.
\end{remark}

\subsection{Why four, in any simply-laced system}

Theorem~\ref{thm:bmstops} was proved by counting transitive factorisations in $S_m$,
which is unavailable outside $B_m$ and $D_m$. The bound has a shorter and more
general reason, and it explains where the number four comes from.

\begin{lemma}[Four blocks at most]\label{lem:fourblocks}
Let $\{r_1,\dots,r_m\}$ be an orthogonal set of roots of a simply-laced system,
normalised to $|r_i|^2=2$, and let $a$ be another root. Then
$\sum_i\langle a,r_i\rangle^2=4$, and since the lattice is integral at most
\textbf{four} of the products are nonzero.
\end{lemma}

\begin{proof}
Write $a=\sum c_ir_i$ with $c_i=\langle a,r_i\rangle/2$. Then
$2=|a|^2=\sum c_i^2|r_i|^2=\tfrac12\sum\langle a,r_i\rangle^2$.
\end{proof}

\begin{corollary}\label{cor:fourblocks}
If a reflection group is generated by an orthogonal set of $m$ roots together with
one further root, the number of irreducible blocks satisfies $s\ge m-3$; so it acts
irreducibly only if $m\le4$.
\end{corollary}

The count is sharp and visible in the root systems: with respect to a maximal
orthogonal set, no root of $E_7$ or of $E_8$ meets all the blocks, while in $D_4$
sixteen of the twenty-four roots do. That is why rank four is the last one where an
irreducible threshold word can occur, and it is the same four as in
Theorem~\ref{thm:bmstops} --- there it looked like a fact about signed permutations,
and it is really a fact about $|r|^2=2$. Checked over $\mathbb{Z}$ by
\path{verificadores/cuatro_bloques.py}.

\subsection{The equation $\det(N-A)=0$, in any dimension --- and why it stops}

In vector rank three, $\det(N-A)=0$ becomes $bw_2=cw_3$ in an eigenbasis of
$A$, a bidegree $(1,1)$ equation with coefficient matrix of rank $2$. The general
formula below shows why this condition becomes identically zero in vector rank four.

\begin{lemma}\label{lem:EA}
Let $A\in\PGL_m$, $c=\pm1$, and $N=c\bigl(2w\tilde w/s-I\bigr)$ with $s=\tilde ww$.
Put $C=cI+A$. Then
\[
  \det(N-A)=0\iff \tilde w\,E_A\,w=0,\qquad
  \boxed{\,E_A:=\det(C)\,I-2c\operatorname{adj}(C)\,}
\]
So $\mathcal C_A=\{N:\det(N-A)=0\}$ is a bidegree $(1,1)$ hypersurface in
$\RP^{m-1}\times\RP^{m-1}$, irreducible of dimension $2(m-1)-1$ as soon as
$\rk E_A\ge2$. When $C$ is invertible, $\operatorname{adj}C=\det(C)C^{-1}$, so
$E_A=\det(C)\bigl(I-2cC^{-1}\bigr)$ and $\rk E_A=\rk(A-cI)$.
\end{lemma}

\begin{proof}
$N-A=-C+\tfrac{2c}{s}w\tilde w$, and $\det(C-t\,w\tilde w)=\det C-t\,\tilde
w\operatorname{adj}(C)w$ for any $C$, invertible or not. Multiply by $s$.
\end{proof}

\begin{corollary}
For $m=3$ and $c=+1$ --- the case of $V_5$ --- with $A$ regular semisimple of
eigenvalues $1,\lambda,\lambda^{-1}$, $C=I+A$ is invertible and in that eigenbasis
$E_A$ is proportional to $\operatorname{diag}(0,1,-1)$, so $\tilde wE_Aw=0$ is
$bw_2=cw_3$ and $\rk E_A=2$.  The involution-chart lemma of
This recovers the rank-three equation directly.
\end{corollary}

\begin{proposition}[Why it stops at $m=4$]\label{prop:stops}
For $m=4$ and $c=-1$ --- the case of $V_6$ --- one has $C=A-I$, and $A\in\mathcal P$
forces $\rk(A-I)\le2=m-2$. The adjugate of a matrix of corank $\ge2$ vanishes, and
so does its determinant. Hence
\[
  E_A=0\quad\text{identically,}
\]
and the condition $\det(N-A)=0$ is \textbf{automatic}: it carries no information.
\end{proposition}

\begin{remark}[Where this leaves the threshold in dimension three]\label{rem:whereleft}
Consequently the fibration modelled on the $V_5$ proof does not close. Its stages
would be $A\in\mathcal P$ ($11$), $N\in c\,\mathrm{Inv}$ with no condition by
Proposition~\ref{prop:stops} ($6$), $H_4\in\mathcal D_Y$ ($3$ by
Proposition~\ref{prop:11hyp}), and the two factorisations ($1+1$): total $22$,
whereas the submersive examples have local dimension $21$. This parameter
count therefore cannot describe a neighbourhood of those examples with
independent parameters of the stated dimensions. The missing compatibility
or fibre constraint is \emph{not} $\det(N-A)=0$; identifying it, and controlling
exceptional components, remain open.

So: irreducibility of $V_6$ is \textbf{not} proved, and neither is the threshold
theorem in dimension three. What is proved unconditionally is
Proposition~\ref{prop:rp3}: the component of $V_6$ through the explicit rational
point is contained in $S_6$, so every configuration on \emph{that} component
preserves a quadric. What is new here is the dimension-free
Lemma~\ref{lem:EA}, which explains where $bw_2=cw_3$ came from, and
Proposition~\ref{prop:stops}, which turns ``the mechanism does not seem to
generalise'' into a theorem about why: in $\mathbb{RP}^3$ the matrix $A-I$ is forced
to have corank two, and adjugates of such matrices vanish.
\end{remark}
\section*{Reproducibility and computational disclosure}
Every displayed rational polytope has an accompanying exact-arithmetic certificate
for scalar monodromy, boundedness, genuine facets, strict first impacts, closure and
primitivity. Floating-point exploration is not used in any proof. Research and manuscript development for all three papers used AI assistance from early
July through late August 2026 (approximate). The models used were OpenAI's GPT-6 Astra,
GPT-5.6 Sol, and GPT-5.6 Luna, and Anthropic's Fable 5 and Opus 5, accessed through
ChatGPT Work, ChatGPT Codex, and Claude Code. No API was used. The assistance supported
research, manuscript drafting, development and checking of verification scripts, and revision.
Jorge Lucas Gonz\'alez reviewed the manuscripts and assumes full responsibility; no AI
system is an author.
\section*{Conclusion}
The paper establishes the parity obstruction, the general rank lower bound,
and a collection of exact higher-dimensional polytope constructions.  It also
separates the proved threshold statements from the irreducibility problem that
remains open in projective three-space.  The examples and their exact
certificates are therefore reusable independently of that unresolved step.
\par\medskip
\noindent\textbf{Competing interests:} The author declares none.
\par\medskip
\begin{thebibliography}{99}
\bibitem{opb} M.~Bialy, C.~Fierobe, A.~Glutsyuk, M.~Levi, M.~Plakhov and
S.~Tabachnikov, \emph{Open Problems on Billiards and Geometric Optics},
Arnold Math. J. \textbf{8}(3--4) (2022), 411--422,
\href{https://doi.org/10.1007/s40598-022-00198-y}{doi:10.1007/s40598-022-00198-y}.
\bibitem{bknz} V.~Blumen, K.~Y.~Kim, J.~Nance and V.~Zharnitsky,
\emph{Three-period orbits in billiards on the surfaces of constant curvature},
Int. Math. Res. Not. (2012), 5014--5024,
\href{https://doi.org/10.1093/imrn/rnr228}{doi:10.1093/imrn/rnr228}.
\bibitem{cb-icm} W.~Crawley-Boevey, \emph{Quiver algebras, weighted projective
lines, and the Deligne--Simpson problem}, Proc. ICM Madrid 2006, vol.~II,
117--129.
\bibitem{fie-ex} C.~Fierobe, \emph{Examples of projective billiards with open
sets of periodic orbits}, Discrete Contin. Dyn. Syst. \textbf{44} (2024),
3287--3301, \href{https://doi.org/10.3934/dcds.2024059}{doi:10.3934/dcds.2024059}.
\bibitem{fie-tri} C.~Fierobe, \emph{On projective billiards with open subsets
of triangular orbits}, Israel J. Math. \textbf{267} (2025), 501--548,
\href{https://doi.org/10.1007/s11856-024-02712-z}{doi:10.1007/s11856-024-02712-z};
\href{https://arxiv.org/abs/2005.02012}{arXiv:2005.02012}.
\bibitem{fie-caustics} C.~Fierobe, \emph{Only quadrics have pseudo-caustics ---
on caustics of Riemannian, pseudo-Euclidean and projective billiards in higher
dimensions}, Math. Z. \textbf{312} (2026), Art. 67,
\href{https://doi.org/10.1007/s00209-026-03939-3}{doi:10.1007/s00209-026-03939-3};
\href{https://arxiv.org/abs/2306.06757}{arXiv:2306.06757}.
\bibitem{heidari-honari} F.~Heidari and B.~Honari, \emph{The three harmonic
homologies theorem}, Expo. Math. \textbf{42} (2024), 125544,
\href{https://doi.org/10.1016/j.exmath.2024.125544}{doi:10.1016/j.exmath.2024.125544}.
\bibitem{tab} S.~Tabachnikov, \emph{Introducing projective billiards},
Ergodic Theory Dynam. Systems \textbf{17} (1997), 957--976,
\href{https://doi.org/10.1017/S0143385797086239}{doi:10.1017/S0143385797086239}.
\end{thebibliography}
\end{document}
